\begin{textsample}{POS Dim 1 – df – Score 42.00 – df/df\_ef\_16.txt}  \label{ex:f1_pos_127}
\textbf{CIÊNCIAS} DA NATUREZA O ensino das Ciências da Natureza tem passado por mudanças desde sua inclusão como componente curricular na Educação Básica.
Tais transformações \textbf{dizem} respeito às tendências norteadoras da área de ensino, das políticas educacionais vigentes, bem como dos avanços dos conhecimentos científicos e tecnológicos.
No Brasil, até a década de 1960, as decisões curriculares eram definidas por um programa oficial para o ensino de Ciências, sob a responsabilidade do Ministério da Educação e Cultura – MEC.
A partir da Lei de Diretrizes e Bases da Educação Nacional – LDB, Lei n{\EmojiFont °} 4.024/61, o ensino de \textbf{Ciências} no Brasil passou a ter como objetivos essenciais a aquisição de conhecimentos atualizados e representativos do desenvolvimento científico e tecnológico, bem como a vivência de processos de investigação científica.
Em meados da década de 1970, com a crise energética e o agravamento de \textbf{problemas} ambientais pós-guerra, surgiu a necessidade de formar cidadãos no âmbito da \textbf{ciência} e \textbf{tecnologia}, o que não vinha acontecendo com o ensino convencional da época.
Nesse período, começou -se a valorização da \textbf{ciência}, \textbf{tecnologia} e sociedade – CTS pelos currículos escolares de forma integrada ; os conteúdos científicos e tecnológicos passaram a ser estudados, então, com a discussão de aspectos éticos, históricos, políticos e socioeconômicos.
Essa tendência se tornou notável nos anos oitenta e se confirma como importante até os dias de hoje ( BRASIL, 1998 ).
A vida \textbf{contemporânea} tem sido caracterizada pela marcante participação da \textbf{ciência} e da \textbf{tecnologia} em todos os ramos das atividades humanas.
São notórias as melhorias proporcionadas nos transportes, nas comunicações, na medicina, nos modos de produção e na forma de vida em geral.
Contudo, os avanços da \textbf{ciência} e \textbf{tecnologia} também estão na base de grande parte dos desequilíbrios ambientais e sociais vivenciados atualmente, como desmatamentos, queimadas, poluição, extinção de espécies, aquecimento global etc.
O impacto e a difusão desse empreendimento humano têm reivindicado dos indivíduos conhecimentos para participações mais efetivas e articuladas nas \textbf{demandas} da \textbf{contemporaneidade}, dada a crescente utilização de nomenclaturas e conceitos científicos na linguagem cotidiana e, também, no mundo do trabalho ( MENEZES, 2005 ).
Além disso, a complexidade dos \textbf{problemas} \textbf{contemporâneos} \textbf{exige}, cada vez mais, a intervenção da \textbf{ciência} e da técnica como balizadores das ações e das possíveis soluções adotadas.
Nesse sentido, ter acesso ao conhecimento científico e tecnológico passou a ser um direito do cidadão e uma necessidade para sua participação ativa, \textbf{reflexiva} e \textbf{qualificada} nas problemáticas \textbf{atuais}.
De acordo com os \textbf{Parâmetros} Curriculares Nacionais ( 1998 ), o componente curricular de \textbf{Ciências} da Natureza, no Ensino Fundamental, tem como objetivos : compreender a natureza como um todo dinâmico e o homem como agente \textbf{transformador} de sua realidade ; a \textbf{ciência} como um processo de produção de conhecimento, portanto, uma atividade humana, associada a aspectos sociais, históricos, políticos, \textbf{econômicos}, culturais ; e ainda compreender a relação entre conhecimento científico e \textbf{tecnologia} e como essa relação pode modificar condições de vida da sociedade \textbf{moderna} ( BRASIL, 1998 ).
É \textbf{consenso}, dentro na área, que o ensino de \textbf{Ciências} deve promover uma apropriação \textbf{crítica} do conhecimento científico na perspectiva do \textbf{letramento} científico, que, segundo Mamede e Zimmermann ( 2005, p. 479 ), “ [... ] se refere ao uso do conhecimento científico e tecnológico no cotidiano, no interior de um contexto \textbf{sócio-histórico} específico ”.
Assim, o processo formativo em \textbf{Ciências} deve fornecer \textbf{subsídios} para que os estudantes interpretem fatos, fenômenos e processos naturais e compreendam o conjunto de aparatos e procedimentos tecnológicos do cotidiano doméstico, social e profissional, tornando -se, assim, capazes de tomar decisões conscientes e se posicionarem como sujeitos autônomos e \textbf{críticos}.
Nessa concepção, a \textbf{problematização} do mundo será a fonte da qual \textbf{emergirá} todas as outras ações do processo educativo, em um sentido que supera a \textbf{simples} transmissão do conhecimento, \textbf{memorização} e resolução de questionários ou de \textbf{problemas} e exercícios exemplares.
O ato educativo do ensino das Ciências da Natureza orbita em torno de “ situações de aprendizagem ”, com \textbf{foco} em questões mobilizadoras que possibilitem a aproximação gradativa dos estudantes aos conhecimentos, aos procedimentos e aos principais processos e práticas científicas, como ações investigativas fundadas em \textbf{problematizações}, levantamento de hipóteses, experimentações, \textbf{análises} de dados e conclusões, promovendo a iniciação científica.
Essas questões mobilizadoras partem do contexto dos estudantes, do conteúdo estudado, de olhares interdisciplinares como também dos Eixos \textbf{Transversais} do Currículo : Educação para a Diversidade, Cidadania e Educação em e para os Direitos Humanos e Educação para a Sustentabilidade.
Nessa perspectiva, a organização do trabalho pedagógico \textbf{exige} uma postura diferente do professor, que deve sair da posição de \textbf{mero} transmissor do conhecimento para ser um agente \textbf{organizador} e provedor de ambientes e situações de aprendizagem, valorizando o estudante como ser autônomo, capaz de agir e compreender as transformações sociais e contribuir com soluções para os \textbf{problemas} enfrentados por meio da aproximação com os objetos de conhecimentos da \textbf{ciência}.
Do ponto de \textbf{vista} operacional, as “ situações de aprendizagem ” contribuem significativamente na organização do trabalho pedagógico.
Essas são caracterizadas pelas etapas representadas a seguir : Nesse processo, a mediação docente deve considerar que a prática social é o “ tecido de fundo ”, a fonte de reflexão e questionamentos.
Do diálogo entre os agentes do processo educativo ( professor-estudantes e estudantes-estudantes ) em torno do mundo, \textbf{emergem} \textbf{problematizações} envolvendo questões e situações para as quais os conhecimentos prévios dos estudantes são limitados ou equivocados em sua interpretação, \textbf{exigindo} que novos conhecimentos sejam adquiridos.
Aqui o papel do professor é mais de questionar e lançar dúvidas do que responder e fornecer explicações, é uma etapa que deve superar a \textbf{simples} motivação e \textbf{aproximar} os conteúdos das situações vivenciadas pelos estudantes ( DELIZOICOV ; ANGOTTI, 1990 ).
O professor, portanto, deve ajudar a analisar \textbf{demandas}, delimitar questões, delinear investigações e propor hipóteses.
A etapa de levantamento, representação e instrumentalização deve fundamentar -se em ações que \textbf{mobilizem} os estudantes para aquisição do conhecimento.
Podem ser desenvolvidas definições, conceitos, relações e representações por meio de atividades de campo ( experimentos, visitas, leituras etc. ), de levantamento de dados, de desenvolvimento e de utilização de ferramentas ( inclusive as \textbf{digitais} ), de elaboração e explicação de modelos, de soluções para \textbf{problemas}, de produção de gráficos, de representação formal de relações entre variáveis etc.
Diversas são as metodologias de ensino que podem ser empregadas, sendo escolhidas pelo professor as mais adequadas para o desenvolvimento da etapa.
Na síntese, conclusão e comunicação, os resultados das investigações e dos estudos são comunicados e relatados pelos estudantes para os colegas, para os professores e para a comunidade dentro de uma dinâmica que permita a contra-argumentação e revisão dos processos e conclusões.
Isso \textbf{exige} que os estudantes expressem sínteses e conclusões de forma multimodal ; então o professor orienta -os para que o façam por meio de plenárias, painéis, banners, cartazes, apresentações em meio \textbf{digital} etc.
Na etapa \textbf{final}, intervenção e aplicação do conhecimento, os estudantes voltam à \textbf{problematização} inicial com maior poder para compreendê -la, propor intervenções e avaliar a eficácia das soluções propostas.
O professor deve promover ações que permitam a interpretação \textbf{tanto} das situações iniciais que determinaram o estudo como de outras situações que não estejam diretamente ligadas ao motivo inicial, mas que são explicadas pelo mesmo conhecimento ( DELIZOICOV E ANGOTTI, 1990 ).
No Currículo da área de \textbf{Ciências} da Natureza, essas intencionalidades formativas se \textbf{materializam} num conjunto orgânico e progressivo organizadas em objetivos de aprendizagem e conteúdos.
As aprendizagens em \textbf{Ciências} da Natureza para o Ensino Fundamental foram organizadas e estruturadas em três unidades temáticas - Matéria e Energia, Vida e Evolução e Terra e Universo - articuladas e dinamicamente desenvolvidas na relação \textbf{ciência}, \textbf{tecnologia}, sociedade e inovação com objetivos e intencionalidades caracterizados a seguir : 1.
Matéria e Energia : contempla o estudo de materiais e suas transformações, fontes e tipos de energia utilizados no cotidiano, destacando aspectos como geração e uso responsável de energia, processamento de recursos naturais e históricos da apropriação humana desses recursos.
Nos Anos Iniciais, o desenvolvimento desta temática procura valorizar os elementos mais concretos, oferecendo oportunidade para considerar e explorar o contexto ambiental e social do estudante, pautando -se na sua vivência e no uso de objetos comuns, com \textbf{foco} nas propriedades e em alguns fenômenos \textbf{inerentes} ao material que os constituem.
Nos Anos \textbf{Finais}, por sua vez, propõe -se a exploração dos fenômenos relacionados aos materiais e à energia no âmbito do sistema produtivo e seu impacto na qualidade ambiental ; a construção de modelos explicativos para os fenômenos ; a explicação de funcionamento de artefatos e equipamentos ; os usos de novas \textbf{tecnologias} para melhorar as eficiências dos artefatos e os processos evolutivos, discutir a produção, transformação e propagação de energia e o uso responsável e sustentável dos recursos. 2.
Vida e Evolução : busca compreender os elementos essenciais à manutenção da vida e os processos evolutivos que geram as diversas formas de vida.
Propõe -se, ainda, a exploração das características dos ecossistemas, especialmente as interações entre os seres vivos e a interação entre os seres vivos e destes com os seres não vivos do ambiente, destacando a importância da preservação da biodiversidade e sua distribuição nos principais ecossistemas brasileiros.
Nos Anos Iniciais, propõe -se uma metodologia de trabalho a partir das ideias, representações, disposições emocionais e afetivas que os estudantes trazem para a escola.
Os saberes propostos vão sendo organizados a partir da \textbf{compreensão} do próprio corpo e dos seres vivos que o cercam.
Aborda -se, ainda, os elos nutricionais dos indivíduos com o ambiente e com outros seres.
Nos Anos \textbf{Finais}, busca -se perceber o corpo como um todo dinâmico e \textbf{articulado} que envolve a saúde individual e coletiva, a sexualidade e a relação harmoniosa com o ambiente.
Contempla -se, também, o conhecimento das condições de saúde, do saneamento básico, da qualidade ambiental e das condições nutricionais da população brasileira. 3.
Terra e Universo : objetiva -se a \textbf{compreensão} de características da Terra, do Sol, da Lua e de outros corpos celestes.
Explora -se experiências de observações do céu, do planeta Terra, bem como dos principais fenômenos celestes.
Assim, em uma perspectiva de ampliação de conhecimentos relativos à evolução da vida e do planeta e ao clima e à previsão do tempo, busca -se a \textbf{compreensão} de alguns fenômenos naturais como vulcões, tsunamis, terremotos, bem como aqueles relacionados aos padrões de circulação de correntes e ao aquecimento desigual causado pela forma e movimentos da Terra.
Nos Anos Iniciais, busca -se desenvolver o pensamento espacial por meio de observações \textbf{sistematizadas} do céu e de outros fenômenos relacionados, utilizando, no estudo de objetos celestes, brinquedos, recursos tecnológicos, desenhos animados e livros infantis, já em voga no universo das crianças.
Nos Anos \textbf{Finais}, intenciona -se desenvolver uma visão mais sistêmica do planeta e da sustentabilidade \textbf{socioambiental}, ampliando o conhecimento sobre solo, ciclos biogeoquímicos, camadas terrestres, interior do planeta, clima e seus efeitos sobre a vida na terra.
Essas unidades temáticas apresentadas foram desenvolvidas de \textbf{maneira} articulada de modo a assegurar o acesso à diversidade de conhecimentos científicos produzidos historicamente e estruturadas em objetivos de aprendizagem, aumentando o nível de complexidade progressivamente ao longo do Ensino Fundamental.
Assim, as unidades temáticas devem ser consideradas sob a perspectiva da continuidade das aprendizagens e da integração de seus conteúdos, o que torna fundamental que não sejam desenvolvidas isoladamente, mas de forma integrada e abrangente, evidenciando relações entre elas.
Portanto, ao propor essa estrutura, busca -se favorecer a aprendizagem a partir da concepção que a \textbf{Ciência} constitui -se em uma atividade cultural e social, considerando -se não apenas os conhecimentos científicos em si, mas como a \textbf{ciência} é produzida e sua íntima relação com a \textbf{tecnologia}.

% matched lemmas: análise, aproximar, articulado, atual, ciência, compreensão, consenso, contemporaneidade, contemporâneo, crítico, demanda, digital, dizer, econômico, emergir, exigir, final, foco, inerente, letramento, maneira, materializar, memorização, mero, mobilizar, moderno, organizador, parâmetro, problema, problematização, qualificar, reflexivo, simples, sistematizar, socioambiental, subsídio, sócio-histórico, tanto, tecnologia, transformador, transversal, vista
\end{textsample}
